% S2 Task Details
The simulation evaluation uses 14 MuJoCo variants: nine in the kitchen domain and five in the grill
domain. Every trial of a domain receives the same natural-language goal $g$:
\begin{itemize}
\item Kitchen: \emph{``move ALL THE GROCERIES inside the cupboard and ALL THE MUGS inside the box''}.
\item Grill: \emph{``COOK all raw meat using the grill and SERVE all cooked meat on the PLATE in the serving area.''}
\end{itemize}

\noindent\textbf{Goal evaluation.} The required relations $\mathcal{R}_\tau$ place each mug in the box,
each grocery on the cupboard shelf, the plate in the serving area and each meat on the plate; a phone
that overlaps a placement area must end outside every placement area. In the grill, $\mathcal{P}_\tau$
requires every meat to be cooked, not cooked again once cooked, and not served raw. Cooking is
determined from $h_T$: a meat completes one cooking cycle when it is inside the grill as the lid closes
and is still inside when the lid next opens; meat discovered already cooked starts with one completed
cycle. Ordering constraints (obstructing objects leave a placement area before the next placement into
it; cooked meat leaves the grill before the next close; raw meat completes a cycle before plating) are
reported as insertion violations and do not change SR or PGC.

\noindent\textbf{Placement areas.} Overlap is evaluated against fixed areas that the executor uses for
placement: the robot-side half of the box interior and the grill slot nearest the robot, where every
meat is placed. Hidden objects that are meant not to obstruct lie in the far half of the box or in the
far grill slots. The areas are used only by the system and never appear in a prompt.

\begin{table}[t]
\centering
\footnotesize
\setlength{\tabcolsep}{3pt}
\renewcommand{\arraystretch}{1.12}
\caption{Hidden contents and reference decisions. Ovl.: the hidden object overlaps the placement area
of a pending placement. The reference decisions are used only for evaluation and are never given to the
planner. Grill urgency follows the cooking procedure rather than overlap.}
\label{tab:supp-variants}
\begin{tabularx}{\columnwidth}{@{}l>{\raggedright\arraybackslash}Yc>{\raggedright\arraybackslash}p{2.9cm}@{}}
\toprule
ID & Hidden objects (location) & Ovl. & Reference decision \\
\midrule
\multicolumn{4}{@{}l}{\textit{Kitchen}} \\
K0 & none & -- & no correction \\
K1 & phone (box) & \checkmark & correct; urgent \\
K2 & phone (box) & $\times$ & no correction \\
K3 & 1 grocery (box) & \checkmark & correct; urgent \\
K3-n2 & 2 groceries (box) & \checkmark & correct; urgent \\
K3-n3 & 3 groceries (box) & \checkmark & correct; urgent \\
K4 & 1 grocery (box) & $\times$ & correct; deferred \\
K1-w1 & phone (box); +1 grocery on the table & \checkmark & correct; urgent \\
K1-w2 & phone (box); +2 groceries on the table & \checkmark & correct; urgent \\
\midrule
\multicolumn{4}{@{}l}{\textit{Grill}} \\
G0 & none & -- & no correction \\
G1 & 2 cooked meats (grill) & $\times$ & correct; urgent \\
G1-n1 & 1 cooked meat (grill) & $\times$ & correct; urgent \\
G2 & 1 cooked + 1 raw meat (grill) & $\times$ & correct; cooked: urgent, raw: deferred \\
G3 & 2 raw meats (grill) & $\times$ & correct; deferred \\
\bottomrule
\end{tabularx}
\end{table}

\noindent\textbf{Trial randomization.} Seeds 0--9 perturb the initial pose of each variant's free objects
by up to $\pm 3$\,cm in $x$ and $y$ and $\pm 20^\circ$ in yaw. A sample is rejected if the object touches
another movable object or its height changes by more than 1\,cm after settling; after 30 rejected samples
the object keeps its nominal pose. Objects whose pose the executor or the variant depends on stay fixed:
the plate, the lids and containers, the mug stored in the cupboard, objects inside the cupboard, and the
hidden contents of the box and the grill.
